%% RESUMO NA LINGUA VERNACULA (elemento obrigatorio -- NBR 6028)
%% Versao dummy: resumo do main.tex da dissertacao real.
%% As palavras-chave sao definidas em config/dados.tex.

\begin{resumo}
O mercado de assessoria de investimentos no Brasil apresenta um duplo gargalo
de distribuição: ferramentas automatizadas tradicionais limitam-se a
questionários estáticos de \textit{suitability}, enquanto modelos de
Inteligência Artificial (IA) Generativa pura apresentam risco de alucinação
matemática incompatível com um ambiente regulado. Nesse contexto, bancos
públicos brasileiros enfrentam um gargalo de capacidade estatal para escalar a
assessoria financeira qualificada ao seu público de perfil popular. O tema
desta pesquisa é uma arquitetura de IA Generativa com recuperação de
conhecimento (\textit{Retrieval-Augmented Generation}, RAG) e cálculo
determinístico para democratizar a assessoria em renda fixa em bancos
públicos. A pergunta de pesquisa é se uma arquitetura de agente único que, por
meio de \textit{function calling}, separa o cálculo quantitativo em Python da
geração de linguagem por \textit{Large Language Model} (LLM) e recupera
normativos vigentes de uma base RAG é capaz de oferecer orientação em renda
fixa a clientes pessoa física de bancos públicos com precisão numérica
verificável e conformidade à Resolução CVM nº~30/2021. A pesquisa adota
abordagem multimétodos, apoiada na metodologia \textit{Design Science
Research}, combinando mapeamento qualitativo de demandas do investidor de
varejo, validação técnica dos módulos de cálculo e da base de conhecimento
regulatória, \textit{backtesting} histórico de 12 meses e avaliação técnica
automatizada por comparação com padrões de referência. Espera-se demonstrar que
a arquitetura proposta é capaz de produzir orientações tecnicamente adequadas,
contribuindo para a literatura de IA aplicada à Administração Pública e para a
democratização do acesso à assessoria financeira no varejo bancário público.
\end{resumo}
